\subsection{线性微分方程积分的线性性}

求解线性微分方程 (2) 是一个线性问题（《代数》，II，第 240 页）；齐次线性方程

\[
{\frac {d \mathbf {x}}{d t}} = A (t). \mathbf {x}\tag{4}
\]

称为与非齐次方程 (2) 相伴的方程；并且知道（《代数》，II，第 241 页，命题 14），若 \(u_{1}\) 是非齐次方程 (2) 的积分，则此方程的每个积分都具有形式 \(u + u_{1}\)，其中 \(u_{1}\) 是相伴齐次方程 (4) 的解，反之亦然。我们将在本小节中首先研究齐次方程 (4) 的积分。

命题 1. 定义在 J 上的齐次线性方程 (4) 的积分所成之集 I 是由 J 到 E 中的连续映射所组成的空间 \(\mathcal{C}(\mathrm{J};\mathrm{E})\) 的一个向量子空间。

证明是直接的。

定理 2. 对 \((t_0, \mathbf{x}_0)\) 的每个点 \(\mathbf{J} \times \mathbf{E}\)，设 \(\mathbf{u}(t, t_0, \mathbf{x}_0)\) 是齐次方程 (4) 的定义在 \(\mathbf{J}\) 上且等于 \(\mathbf{x}_0\) 于 \(t_0\) 的积分。

\(1^{\circ}\) 对每个点 \(t \in \mathbf{J}\)，映射 \(\mathbf{x}_0 \mapsto \mathbf{u}(t, t_0, \mathbf{x}_0)\) 是从 \(C(t, t_0)\) 到其自身的双射双连续线性映射 \(\mathbf{E}\)。

\(2^{\circ}\) 映射 \(t \mapsto C(t, t_0)\)（从 \(\mathbf{J}\) 到 \(\mathcal{L}(\mathbf{E})\) 中）等同于齐次线性微分方程的积分

\[
\frac {d U}{d t} = A (t) U\tag{5}
\]

它取值 \(I\)（即 \(E\) 到其自身的恒等映射）于点 \(t_0\) 。

\(3^{\circ}\) 对任意点 \(s, t, u\)（属于 \(\mathbf{J}\)）

\[
C (s, u) = C (s, t) C (t, u), \quad C (s, t) = (C (t, s)) ^ {- 1}.\tag{6}
\]

由命题 1，\(\mathbf{u}(t,t_0,\mathbf{x}_1) + \mathbf{u}(t,t_0,\mathbf{x}_2)\)（相应地，\(\lambda \mathbf{u}(t,t_0,\mathbf{x}_0)\) ）是 (4) 的积分且取值 \(\mathbf{x}_1 + \mathbf{x}_2\)（相应地，\(\lambda \mathbf{x}_0\) ）于 \(t_0\)，故由 IV 第 179 页的定理 1，它与 \(\mathbf{u}(t,t_0,\mathbf{x}_1 + \mathbf{x}_2)\)（相应地，\(\mathbf{u}(t,t_0,\lambda \mathbf{x}_0)\) ）相同；于是映射 \(\mathbf{x}_0\mapsto \mathbf{u}(t,t_0,\mathbf{x}_0)\) 是 E 到其自身中的线性映射 \(C(t,t_0)\)，并且可以写 \(\mathbf{u}(t,t_0,\mathbf{x}_0) = C(t,t_0).\mathbf{x}_0\) 。

由于映射 \((X, Y) \mapsto XY\)（从 \(\mathcal{L}(\mathrm{E}) \times \mathcal{L}(\mathrm{E})\) 到 \(\mathcal{L}(\mathrm{E})\) 中）是连续的（《一般拓扑学》，X，第 298 页，命题 8），映射 \(t \mapsto A(t)U\)（从 J 到 \(\mathcal{L}(\mathrm{E})\) 中）对一切 \(U \in \mathcal{L}(\mathrm{E})\) 都是规则的；此外（《一般拓扑学》，X，第 296 页）

\[
\| A (t) X - A (t) Y \| = \| A (t) (X - Y) \| \leqslant \| A (t) \| \| X - Y \|,
\]

于是可以把 IV 第 179 页的定理 1 应用于齐次线性方程 (5)；设 \(V(t)\) 是此方程的定义在 J 上且等于 \(I\) 于 \(t_0\) 的积分。我们有（I，第 6 页，命题 3）

\[
\frac {d}{d t} \big (V (t) \mathbf {x} _ {0} \big) = \frac {d V (t)}{d t} \mathbf {x} _ {0} = A (t) \big (V (t) \mathbf {x} _ {0} \big)
\]

而对 \(t = t_{0}\) 有 \(V(t)\mathbf{x}_{0} = I\mathbf{x}_{0} = \mathbf{x}_{0}\) ；由 IV 第 179 页的定理 1，必有 \(V(t)\mathbf{x}_{0} = C(t, t_{0})\mathbf{x}_{0}\) 对一切 \(x_{0} \in E\)，即 \(V(t) = C(t, t_{0})\) ；这就证明了 \(C(t, t_{0})\) 属于 \(\mathcal{L}(\mathrm{E})\)，换言之，\(x_{0} \mapsto C(t, t_{0})\mathbf{x}_{0}\) 在 E 上连续，并且映射 \(t \mapsto C(t, t_{0})\) 是 (5) 的、等于 I 于 \(t_{0}\) 的积分。

最后，(4) 的积分 \(s \mapsto C(s, u). \mathbf{x}_0\) 等于 \(C(t, u). \mathbf{x}_0\) 于点 \(t\) ，于是由定义，

\[
C (s, u). \mathbf {x} _ {0} = C (s, t) \bigl (C (t, u). \mathbf {x} _ {0} \bigr) = \bigl (C (s, t)   C (t, u) \bigr). \mathbf {x} _ {0}
\]

对任何 \(x_{0} \in E\) 都成立，由此得第一个关系式 (6)；由于 \(C(s, s) = I\)，有 \(C(s, t)C(t, s) = I\) （对 J 中任何 s 与 t）；这就证明了（《集合论》，II，第 86 页，推论）\(C(t, t_{0})\) 是 E 到其自身上的双射映射，其逆映射为 \(C(t_{0}, t)\)。这就完成了定理的证明。

我们称 \(C(t, t_0)\) 为 IV 第 178 页方程 (2) 的预解式。

推论 1. 使每个点 \(\mathbf{x}_0 \in E\) 对应于定义在 J 上的连续函数 \(t \mapsto C(t, t_0). \mathbf{x}_0\) 的映射，是赋范空间 \(E\) 到 (4) 的积分所成的向量空间 \(\mathcal{I}\) （赋以紧收敛拓扑）上的同构。

它当然是 E 到 \(\mathcal{I}\) 上的双射线性映射；而 \(C(t, t_0)\) 在紧集 \(\mathrm{K} \subset \mathrm{J}\) 上有界，故对任何 \(\|C(t, t_0). \mathbf{x}_0\| \leqslant \mathrm{M} \| \mathbf{x}_0\|\) 与 \(t \in \mathrm{K}\) 有 \(\mathbf{x}_0 \in \mathrm{E}\)，这表明此映射是连续的；又由于

\[
C (t _ {0}, t _ {0}). \mathbf {x} _ {0} = \mathbf {x} _ {0},
\]

显然逆映射也是连续的。

推论 2. 映射 \((s,t)\mapsto C(s,t)\)（从 \(\mathrm{J}\times \mathrm{J}\) 到 \(\mathcal{L}(\mathrm{E})\) 中）是连续的。

由 (6) 有 \(C(s, t) = C(s, t_0) \left( C(t, t_0) \right)^{-1}\) ；而映射 \((X, Y) \mapsto XY\)（从 \(\mathcal{L}(\mathrm{E}) \times \mathcal{L}(\mathrm{E})\) 到 \(\mathcal{L}(\mathrm{E})\) 中）是连续的，映射 \(X \mapsto X^{-1}\)（由 \(\mathcal{L}(\mathrm{E})\) 的可逆元素所成的（开）群到其自身上）也是连续的（《一般拓扑学》，IX，第 40 页，命题 14）。

我们可以注意到，映射

\[
t \mapsto C (t _ {0}, t) = (C (t, t _ {0})) ^ {- 1}
\]

有等于 \(-(C(t,t_{0}))^{-1}(dC(t,t_{0})/dt)(C(t,t_{0}))^{-1}\) 的导数（在一个可数集的补集上）（I，第 8 页，命题 4），也就是说（由 IV 第 179 页公式 (5)）等于 \(-C(t_{0},t)A(t)\) 。

推论 3. 设 K 是包含于 J 中的紧区间，并设 \(k = \sup_{t \in \mathbf{K}} \| A(t) \|\) 。

对所有 \(t\) 与 \(t_0\)（都属于 \(\mathbf{K}\)）

\[
\left\| C \left(t, t _ {0}\right) - I \right\| \leqslant e ^ {k \left| t - t _ {0} \right|} - 1.\tag{7}
\]

的确，\(\|A(t)\mathbf{x}_{0}\|\leqslant k\|\mathbf{x}_{0}\|\) 对一切 \(t\in K\) 成立；于是在 K 上，等于 \(x_{0}\) 的常函数便是精确到 \(k\|\mathbf{x}_{0}\|\) 的近似积分（依据 IV 第 18 页的方程 (4)）；由 IV 第 170 页的公式 (15)，便有

\[
\left\| C \left(t, t _ {0}\right) \mathbf {x} _ {0} - \mathbf {x} _ {0} \right\| \leqslant \left\| \mathbf {x} _ {0} \right\| \left(e ^ {k | t - t _ {0} |} - 1\right)
\]

对 K 中的任何 t 与 \(t_{0}\)，以及 E 中的 \(x_{0}\)，都成立；由 \(\mathcal{L}(\mathrm{E})\) 上范数的定义，这等价于不等式 (7)。

命题 2. 设 B 是 E 的连续自同态，它与 t 无关，并且与 \(A(t)\) 对一切 \(t \in J\) 交换；那么 B 与 \(C(t, t_{0})\) 对 J 中一切 t 与 \(t_{0}\) 交换。

的确，由 (5)

\[
\frac {d}{d t} (B C) = B A C = A B C \quad \text { and } \quad \frac {d}{d t} (C B) = A C B,
\]

于是 \(\frac{d}{dt}\bigl (BC - CB\bigr) = A(BC - CB)\) ；但 \(BC(t_0,t_0) - C(t_0,t_0)B = 0\) ，故（IV，第 179 页，定理 1）\(BC(t,t_0) - C(t,t_0)B = 0\) 对一切 \(t\in \mathrm{J}\) 成立。

命题 2 的一个重要特例是：E 被赋予关于复数域 C 的赋范向量空间结构，并且对每个 \(t \in J\)，\(A(t)\) 是 E 关于该向量空间结构的自同态；这就意味着 \(A(t)\) 与 E 的连续自同态 \(x \mapsto i x\) （关于 R 上的向量空间结构）交换；于是 \(C(t, t_{0})\) 与该自同态交换，而这意味着对 J 中任何 t 与 \(t_{0}\)，映射 \(C(t, t_{0})\) 都是 E 关于 C 上赋范向量空间结构的连续自同态。

